\documentclass[../../main.tex]{subfiles}

\begin{document}

\chapter{Evaluasi, Refleksi, dan Integrasi Kompetensi}
\label{chap:evaluasi-akhir}

% =========================================================
% OUTCOME PEMBELAJARAN KHUSUS BAB
% =========================================================

\begin{subcpmkbox}
Bab ini mengintegrasikan seluruh Sub-CPMK, yaitu \SubCPMK{1.1}
sampai \SubCPMK{4.3}, dan menyediakan instrumen untuk mengukur
ketercapaian kompetensi secara menyeluruh.

\vspace{0.4em}
\textbf{CPMK yang didukung:} \CPMK{01}, \CPMK{02}, \CPMK{03},
\CPMK{04} \quad\textbullet\quad
\textbf{CPL yang didukung:} \CPL{02}, \CPL{03}
\end{subcpmkbox}

\begin{tujuanbox}
Setelah mempelajari bab ini, mahasiswa mampu:
\begin{enumerate}
    \item memetakan kembali perjalanan kompetensi dari \CPL{02} dan
          \CPL{03} sampai ke setiap Sub-CPMK;
    \item mempersiapkan diri menghadapi Ujian Tengah Semester dan
          Ujian Akhir Semester berdasarkan kisi-kisi yang tersedia;
    \item merencanakan dan melaksanakan proyek akhir pemrograman
          assembly sesuai rubrik yang ditetapkan;
    \item melakukan penilaian diri (\eng{self-assessment}) dan
          penilaian antarteman (\eng{peer assessment}) secara jujur
          dan berbasis bukti;
    \item menyusun portofolio dan jurnal refleksi yang lengkap;
    \item merencanakan tindak lanjut belajar pada topik lanjutan.
\end{enumerate}
\end{tujuanbox}

\begin{materibox}
\begin{enumerate}
    \item Rekapitulasi capaian pembelajaran dan peta ketercapaian
          kompetensi.
    \item Kisi-kisi UTS, kisi-kisi UAS, dan pedoman proyek akhir.
    \item Rubrik penilaian komprehensif berbasis OBE.
    \item Penilaian diri, penilaian antarteman, portofolio, dan
          jurnal refleksi.
    \item Rekomendasi tindak lanjut dan topik lanjutan.
\end{enumerate}
\end{materibox}

\subfile{section-01}
\subfile{section-02}
\subfile{section-03}

\section{Contoh Penilaian Proyek Akhir}
\label{sec:contoh-bab17}

Agar rubrik pada Tabel~\ref{tab:rubrik-proyek} tidak menjadi
dokumen yang abstrak, berikut contoh penerapannya pada satu proyek.

\subsection{Deskripsi Proyek yang Dinilai}

Seorang mahasiswa mengerjakan Pilihan 1, yaitu kalkulator dengan
validasi masukan. Proyeknya memuat:

\begin{itemize}
    \item prosedur \code{BacaAngka} yang membaca string masukan
          dengan \code{INT 21h} fungsi \hex{0A} lalu mengubahnya
          menjadi nilai biner;
    \item prosedur \code{Hitung} yang menerima dua parameter melalui
          stack beserta kode operator, lalu mengembalikan hasil pada
          \reg{AX} dan \reg{DX};
    \item prosedur \code{TulisAngka} yang mengubah hasil biner
          kembali menjadi deretan digit dan mencetaknya;
    \item penanganan pembagian dengan nol yang mencetak pesan
          kesalahan dan kembali ke menu.
\end{itemize}

\subsection{Penerapan Rubrik}

Tabel~\ref{tab:contoh-penilaian} memperlihatkan cara pemberian skor
beserta alasan yang harus dituliskan penilai. Kolom alasan inilah
yang membuat penilaian dapat dipertanggungjawabkan.

\begin{table}[H]
    \centering
    \caption{Contoh penerapan rubrik proyek akhir.}
    \label{tab:contoh-penilaian}
    \footnotesize
    \begin{tabular}{@{}p{2.9cm}cp{7.6cm}@{}}
        \toprule
        \textbf{Kriteria} & \textbf{Skor} & \textbf{Alasan} \\
        \midrule
        Kebenaran fungsi & 4 &
        Empat operasi berjalan benar; pembagian dengan nol ditangani
        dengan pesan kesalahan; masukan bukan angka ditolak \\
        \addlinespace
        Kualitas kode & 3 &
        Prosedur terpisah dengan baik, tetapi prosedur
        \code{TulisAngka} memiliki tiga baris yang berulang dan tidak
        dijadikan subprosedur \\
        \addlinespace
        Penggunaan instruksi & 4 &
        Perkalian pangkat dua memakai \asm{SHL} alih-alih
        \asm{MUL}; mode \eng{base-relative} dipakai untuk menelusuri
        array digit \\
        \addlinespace
        Pengujian dan \eng{debugging} & 3 &
        Tabel \eng{trace} disertakan untuk dua dari tiga bug; satu
        bug terakhir ditemukan dengan bantuan dosen \\
        \addlinespace
        Laporan dan presentasi & 4 &
        Laporan memuat rancangan, kendala, dan bukti pengujian;
        presentasi menjelaskan alasan pemilihan mode pengalamatan \\
        \midrule
        \textbf{Rata-rata} & \textbf{3{,}6} &
        Setara dengan kategori ``baik'' menuju ``sangat baik'' \\
        \bottomrule
    \end{tabular}
\end{table}

\subsection{Pelajaran dari Contoh}

Tiga hal dapat dipelajari dari contoh di atas.

\begin{enumerate}
    \item \textbf{Skor tertinggi tidak menuntut proyek yang paling
          rumit.} Proyek sederhana yang dikerjakan dengan benar dan
          didokumentasikan dengan baik memperoleh nilai lebih tinggi
          daripada proyek rumit yang tidak dapat dijelaskan.

    \item \textbf{Alasan harus dituliskan.} Rubrik tanpa alasan
          hanya menghasilkan angka; rubrik dengan alasan
          menghasilkan umpan balik yang dapat ditindaklanjuti.

    \item \textbf{Bukti menentukan.} Kriteria pengujian dan
          \eng{debugging} tidak dapat dinilai tanpa tabel
          \eng{trace} dan catatan bug. Inilah alasan jurnal
          praktikum wajib dipelihara sejak pertemuan pertama.
\end{enumerate}

\begin{infobox}
Perhatikan bahwa seluruh kriteria pada rubrik berasal dari CPMK.
Kriteria ``penggunaan instruksi'' misalnya berasal dari \CPMK{02},
sedangkan ``kebenaran fungsi'' dan ``kualitas kode'' berasal dari
\CPMK{03}. Tidak ada kriteria yang tidak dapat ditelusuri ke CPMK
--- inilah yang dimaksud dengan asesmen berbasis OBE.
\end{infobox}

\section{Pelaksanaan Proyek Akhir}
\label{sec:praktikum-bab17}

\begin{praktikumbox}
\vspace{1.1em}
\textbf{Judul:} Proyek akhir pemrograman assembly 8086.\\[0.3em]
\textbf{CPMK:} \CPMK{03} dan \CPMK{04} \quad\textbullet\quad
\textbf{Alokasi waktu:} 3 minggu (Minggu 13--15) \quad\textbullet\quad
\textbf{Bentuk:} individu atau berpasangan
\end{praktikumbox}

\subsection{Tujuan}

Menghasilkan satu program assembly yang utuh, benar, teruji, dan
dapat dipertanggungjawabkan sebagai bukti pencapaian CPMK03 dan
CPMK04.

\subsection{Tahapan Pengerjaan}

\begin{enumerate}
    \item \textbf{Minggu 13 --- perencanaan.}
    \begin{enumerate}
        \item Pilih salah satu dari empat pilihan proyek pada
              Bagian~\ref{sec:asesmen-akhir}.
        \item Tulis rancangan satu halaman: tujuan, daftar prosedur,
              masukan/keluaran, dan kasus salah.
        \item Kumpulkan rancangan untuk memperoleh persetujuan.
    \end{enumerate}

    \item \textbf{Minggu 14 --- implementasi dan pengujian.}
    \begin{enumerate}
        \item Implementasikan prosedur satu per satu dan uji segera.
        \item Siapkan minimal dua kasus uji per prosedur: satu benar,
              satu salah.
        \item Catat setiap bug dan cara menemukannya.
        \item Simpan setiap versi \file{.asm} yang berfungsi.
    \end{enumerate}

    \item \textbf{Minggu 15 --- penyelesaian dan presentasi.}
    \begin{enumerate}
        \item Lengkapi komentar kode.
        \item Susun laporan: tujuan, rancangan, kendala, pengujian.
        \item Siapkan demonstrasi maksimal 10 menit.
        \item Tunjukkan satu kasus salah yang ditangani dan jelaskan
              satu keputusan teknis.
    \end{enumerate}
\end{enumerate}

\subsection{Kasus Uji Minimum}

\begin{table}[H]
    \centering
    \caption{Kasus uji minimum yang wajib disiapkan.}
    \label{tab:kasus-uji}
    \small
    \begin{tabular}{@{}p{3.4cm}p{5.2cm}p{4.6cm}@{}}
        \toprule
        \textbf{Jenis kasus} & \textbf{Contoh} &
        \textbf{Perilaku yang diharapkan} \\
        \midrule
        Masukan normal & Nilai dalam rentang wajar &
        Keluaran benar \\
        \addlinespace
        Masukan batas & Nilai terbesar/terkecil &
        Tidak terjadi \eng{overflow} \\
        \addlinespace
        Masukan salah & Huruf pada tempat angka, pembagi nol &
        Ditolak dengan pesan, program tidak berhenti \\
        \addlinespace
        Masukan kosong & Pengguna menekan Enter &
        Tidak membaca sampah penyangga \\
        \bottomrule
    \end{tabular}
\end{table}

\subsection{Hasil yang Dikumpulkan}

\begin{enumerate}
    \item Berkas \file{.asm} proyek yang dapat dirakit tanpa
          kesalahan.
    \item Laporan singkat 2--4 halaman.
    \item Catatan pengujian berisi tabel kasus dan hasil.
    \item Catatan \eng{debugging} berisi bug dan cara penemuan.
    \item Bahan presentasi.
\end{enumerate}

\begin{asesmenbox}
\vspace{1.1em}
\textbf{Asesmen proyek akhir (bobot 15\% dari nilai mata kuliah).}
Penilaian memakai rubrik pada Tabel~\ref{tab:rubrik-proyek}. Contoh
penerapannya ada pada Bagian~\ref{sec:contoh-bab17}.

\begin{tabular}{@{}p{7.4cm}p{5.8cm}@{}}
    \toprule
    \textbf{Kriteria} & \textbf{Bobot terhadap nilai proyek} \\
    \midrule
    Kebenaran fungsi (\CPMK{03}) & 1/5 \\
    \addlinespace
    Kualitas kode (\CPMK{03}) & 1/5 \\
    \addlinespace
    Instruksi dan pengalamatan (\CPMK{02}, \CPMK{03}) & 1/5 \\
    \addlinespace
    Pengujian dan \eng{debugging} (\CPMK{04}) & 1/5 \\
    \addlinespace
    Laporan dan presentasi & 1/5 \\
    \bottomrule
\end{tabular}
\end{asesmenbox}

\section{Latihan}
\label{sec:latihan-bab17}

Latihan pada bab ini bersifat reflektif dan integratif. Tidak ada
jawaban tunggal yang benar; yang dinilai adalah kedalaman dan
kejujuran analisis Anda, serta ketepatan bukti yang ditunjuk.

\begin{enumerate}
    \item \textbf{Integrasi konsep.} Jelaskan bagaimana
          Bab~\ref{chap:segmentasi}, Bab~\ref{chap:mode-pengalamatan},
          dan Bab~\ref{chap:prosedur} saling terkait ketika sebuah
          prosedur menerima parameter melalui stack. Urutkan
          konsepnya dan tunjukkan di mana setiap bab berperan.

    \item \textbf{Analisis kode lintas bab.} Diberikan sebuah
          program yang membaca nama pengguna, menghitung panjangnya,
          lalu mencetak salam. Identifikasi bab mana yang menyediakan
          konsep untuk setiap bagian program tersebut.

    \item \textbf{Penilaian diri.} Isi lembar penilaian diri pada
          Tabel~\ref{tab:self-assessment}. Untuk setiap CPMK, sebutkan
          satu bukti konkret berupa nama berkas atau nomor tugas.
          Klaim tanpa bukti tidak dihitung.

    \item \textbf{Analisis kesalahan.} Pilih satu bug yang pernah
          Anda alami sepanjang semester. Jelaskan: gejala yang
          muncul, hipotesis pertama Anda, langkah yang membuktikan
          atau menggugurkan hipotesis itu, dan penyebab sebenarnya.

    \item \textbf{Perbandingan pendekatan.} Ambil satu program yang
          Anda tulis pada Bab~\ref{chap:kontrol}. Tuliskan ulang
          bagian perulangannya dengan pendekatan berbeda (misalnya
          mengganti \asm{LOOP} dengan \asm{DEC} dan \asm{JNZ}), lalu
          bandingkan jumlah instruksi dan keterbacaannya.

    \item \textbf{Evaluasi rubrik.} Baca kembali rubrik proyek pada
          Tabel~\ref{tab:rubrik-proyek}. Untuk setiap kriteria,
          tuliskan satu hal konkret yang akan Anda lakukan pada
          proyek akhir agar memperoleh skor 4.

    \item \textbf{Rencana tindak lanjut.} Berdasarkan
          Tabel~\ref{tab:tindak-lanjut}, pilih satu jalur lanjutan,
          lalu susun rencana belajar tiga bulan pertama: topik,
          sumber, dan bukti yang akan Anda hasilkan.

    \item \textbf{Kritik buku.} Sebutkan satu bagian buku ini yang
          menurut Anda paling kurang jelas, dan usulkan cara
          perbaikannya. Sertakan nomor bab dan bagian.
\end{enumerate}

\begin{refleksibox}
Jurnal refleksi akhir semester. Jawab dengan jujur dan sertakan
bukti.
\begin{enumerate}
    \item Kompetensi yang paling berkembang pada diri saya adalah
          \dots, dengan bukti \dots
    \item Kompetensi yang masih paling lemah adalah \dots, dan
          rencana perbaikan saya adalah \dots
    \item Kebiasaan belajar yang paling membantu saya adalah \dots
    \item Kebiasaan belajar yang paling merugikan saya adalah \dots
    \item Setelah mempelajari level instruksi, hal yang paling
          mengubah cara saya memandang program adalah \dots
\end{enumerate}
\end{refleksibox}

\section{Rangkuman}
\label{sec:rangkuman-bab17}

Pada bab ini telah dibahas:

\begin{itemize}
    \item Perjalanan kompetensi mata kuliah ini bergerak dari
          pemahaman mesin (\CPMK{01}), ke hubungan instruksi dengan
          perangkat keras (\CPMK{02}), ke penulisan program
          (\CPMK{03}), dan berakhir pada kemampuan memperbaiki
          program (\CPMK{04}).

    \item Setiap Sub-CPMK memiliki bukti belajar yang harus
          ditunjukkan. Bukti tersebut bersifat kumulatif dan
          dipelihara sepanjang semester.

    \item Ujian Tengah Semester mengukur Sub-CPMK 1.1 sampai 2.3,
          sedangkan Ujian Akhir Semester mengukur Sub-CPMK 3.1
          sampai 4.3 secara komprehensif dengan penekanan pada
          kemampuan menulis dan memperbaiki program.

    \item Proyek akhir menawarkan empat pilihan dan dinilai dengan
          rubrik lima kriteria yang seluruhnya dapat ditelusuri ke
          CPMK.

    \item Penilaian diri dan penilaian antarteman adalah bagian
          resmi penilaian dan harus disertai bukti.

    \item Nilai rubrik tanpa alasan hanya menghasilkan angka; rubrik
          dengan alasan menghasilkan umpan balik yang dapat
          ditindaklanjuti.

    \item Tindak lanjut dapat diarahkan ke arsitektur komputer,
          pemrograman sistem, analisis keamanan, atau sistem
          tertanam, sesuai minat mahasiswa.
\end{itemize}

\begin{checklistbox}
\textbf{Checklist Kompetensi Akhir Mata Kuliah}

Tandai butir yang benar-benar dapat Anda buktikan dengan berkas
atau catatan.

\begin{itemize}
    \item[\checkbox] Saya dapat menjelaskan arsitektur 8086, register,
          dan segmentasi memori, serta menunjukkan diagram dan
          perhitungan yang mendukungnya.
    \item[\checkbox] Saya dapat membaca format instruksi mesin,
          menentukan mode pengalamatan, dan menjelaskan alur translasi
          program.
    \item[\checkbox] Saya dapat menulis program assembly lengkap yang
          memuat aritmatika, logika, struktur kontrol, dan prosedur
          berparameter.
    \item[\checkbox] Saya dapat menggunakan interupsi BIOS/DOS untuk
          operasi masukan dan keluaran.
    \item[\checkbox] Saya dapat menelusuri program langkah demi
          langkah, membaca isi memori dan stack, serta menemukan
          penyebab kesalahan.
    \item[\checkbox] Saya dapat membaca hasil \eng{disassembly} dan
          mengusulkan perbaikan yang terukur.
    \item[\checkbox] Portofolio, jurnal praktikum, dan laporan proyek
          akhir saya lengkap.
\end{itemize}

\vspace{0.4em}
Seluruh butir bertanda berarti seluruh CPMK mata kuliah ini telah
tercapai.
\end{checklistbox}

\begin{infobox}
\textbf{Rujukan bab.} Untuk memperluas pemahaman mengenai organisasi
komputer yang menjadi dasar mata kuliah ini, lihat
\cite{stallings2016} dan \cite{tanenbaum2013}. Untuk jalur
pemrograman sistem dan analisis level rendah, lihat
\cite{bartlett2004}, \cite{bartlett2023}, dan
\cite{wikibooksdisasm}. Untuk rujukan teknis resmi keluarga x86,
lihat \cite{intel2024}.
\end{infobox}

\section{Evaluasi}
\label{sec:evaluasi-bab17}

Evaluasi pada bab ini berbentuk \textbf{ujian komprehensif}. Soal
menuntut penggabungan konsep dari beberapa bab sekaligus.

\subsection{Soal Komprehensif}

\begin{enumerate}
    \item \textbf{Analisis terpadu.} Diberikan potongan program
          berikut.

\begin{lstlisting}[language=MASM, numbers=none]
.MODEL SMALL
.STACK 100h
.DATA
    data1 DW 0A000h
    data2 DW 0004h
.CODE
START:
    MOV AX, @DATA
    MOV DS, AX
    MOV AX, data1
    MOV BX, data2
    ADD AX, BX
    MOV AH, 4Ch
    INT 21h
END START
\end{lstlisting}

          Jawab pertanyaan berikut.
          \begin{enumerate}
              \item Berapa nilai \reg{AX} setelah \asm{ADD AX, BX}?
              \item Nyatakan nilai tersebut dalam heksadesimal dan
                    tafsiran bertanda.
              \item Tentukan \flag{CF}, \flag{ZF}, \flag{SF}, dan
                    \flag{OF} beserta alasannya.
              \item Bila \code{data2} diubah menjadi \hex{C000},
                    flag mana yang berubah dan mengapa?
          \end{enumerate}

    \item \textbf{Perancangan prosedur.} Rancang prosedur
          \code{RataRata} yang menerima tiga bilangan 16-bit melalui
          stack dan mengembalikan rata-ratanya pada \reg{AX}.
          Tuliskan gambar isi stack, kode prosedur dengan
          \eng{prologue}/\eng{epilogue}, dan cara pemanggilannya.

    \item \textbf{Perbaikan bug.} Program bermaksud menjumlahkan 10
          elemen array, tetapi hasilnya selalu kurang satu elemen.
          Tuliskan langkah sistematis menemukan penyebabnya dan
          perintah \eng{debugger} yang dipakai.

    \item \textbf{Array, string, dan pencetakan.} Diketahui array
          word dan string. Tuliskan rancangan program yang
          menjumlahkan array, menyalin string, lalu mencetak hasil
          penjumlahan dalam desimal. Sebutkan prosedur dan register
          yang dipakai.

    \item \textbf{Interupsi.} Jelaskan urutan penyimpanan dan
          pemulihan \reg{FLAGS}, \reg{CS}, dan \reg{IP} ketika
          \code{INT 21h} dan \code{IRET} dijalankan. Sertakan diagram
          stack.

    \item \textbf{Perancangan.} Anda diminta membuat program yang
          menampilkan isi array dalam bentuk heksadesimal. Tuliskan
          daftar prosedur beserta masukan dan keluarannya sebelum
          menulis kode.
\end{enumerate}

\subsection{Kisi-kisi Ujian Komprehensif}

\begin{table}[H]
    \centering
    \caption{Kisi-kisi ujian komprehensif
    Bab~\ref{chap:evaluasi-akhir}.}
    \label{tab:kisi-bab17}
    \small
    \begin{tabular}{@{}p{4.0cm}p{5.4cm}c@{}}
        \toprule
        \textbf{Aspek yang diukur} & \textbf{Soal} & \textbf{Bobot} \\
        \midrule
        Integrasi aritmatika, representasi, dan \eng{flag} & Soal 1 &
        20\% \\
        \addlinespace
        Perancangan dan penulisan prosedur & Soal 2, 6 & 25\% \\
        \addlinespace
        Metodologi \eng{debugging} & Soal 3 & 15\% \\
        \addlinespace
        Array, string, dan pencetakan bilangan & Soal 4 & 20\% \\
        \addlinespace
        Mekanisme interupsi & Soal 5 & 20\% \\
        \bottomrule
    \end{tabular}
\end{table}

\begin{checklistbox}
\vspace{1.1em}
\textbf{Checklist Kesiapan Menghadapi UAS}

\begin{itemize}
    \item[\checkbox] Saya dapat menulis program aritmatika dan
          manipulasi bit tanpa melihat contoh.
    \item[\checkbox] Saya dapat menerjemahkan pseudocode bercabang
          dan berulang menjadi assembly.
    \item[\checkbox] Saya dapat menggambar isi stack dan menghitung
          parameter melalui \reg{BP}.
    \item[\checkbox] Saya dapat menelusuri program dengan
          \eng{debugger} dan mengisi tabel \eng{trace}.
    \item[\checkbox] Saya dapat mengolah array/string dan mencetak
          bilangan.
    \item[\checkbox] Saya dapat menjelaskan mekanisme interupsi
          beserta perubahannya pada stack.
\end{itemize}
\end{checklistbox}


\end{document}
